% =============================================================================
% Project : THE PHILOSOPHY OF VIABILITY
% Book    : Book I — Ontology of Viability
% Book ID : VS.B01
% Part    : Part I — Identity and Manifestation
% Part ID : VS.B01.P01
% Chapter : Chapter 08 — Comparative Philosophical Examination
% Chapter ID : VS.B01.P01.C08
%
% File :
%   4_ASSETS/diagrams/
%   VS.B01.P01.C08.S10.SS07.P006.DG01.kant_synthesis_and_apperceptive_unity.tex
%
% Section :
%   Section 10 — Immanuel Kant
%
% Asset :
%   Diagram DG01 — Kant's Synthesis and Apperceptive Unity
%
% Structural Identifier :
%   VS.B01.P01.C08.S10.SS07.P006.DG01
%
% Asset Type : DG
%
% Purpose :
%   To distinguish the manifold of intuition, synthesis,
%   the transcendental unity of apperception, and the
%   representation of an object in possible experience.
%
% Intellectual Benefit :
%   The diagram makes visible the structure of Kant's
%   transcendental account and clarifies the relation
%   between apperceptive unity and objectivity.
%
% Aesthetic / Experiential Benefit :
%   It provides a visual pause within a conceptually dense
%   discussion and gives the reader a compact overview of
%   the central argument.
%
% Representational Limit :
%   The diagram represents the conditions of possible
%   experience in Kant's account. It does not depict an
%   ontological production of beings by the subject.
%   The lower distinction explicitly marks the limit
%   of the inference from apperceptive unity to a
%   substantial or numerically persistent soul.
%
% Dependencies :
%   TikZ; positioning library.
%   Both are used elsewhere in the project.
%
% =============================================================================

\begin{figure}[htbp]
    \centering

    \begin{tikzpicture}[
        font=\small,
        node distance=0.68cm,
        process/.style={
            draw=blue!55!black,
            fill=blue!6,
            rounded corners=3pt,
            line width=0.7pt,
            text width=9.0cm,
            minimum height=0.82cm,
            align=center,
            inner sep=6pt
        },
        central/.style={
            process,
            fill=blue!15,
            font=\small\bfseries
        },
        result/.style={
            draw=teal!55!black,
            fill=teal!7,
            rounded corners=3pt,
            line width=0.7pt,
            text width=9.0cm,
            minimum height=0.82cm,
            align=center,
            inner sep=6pt
        },
        limit/.style={
            draw=black!55,
            fill=black!4,
            rounded corners=3pt,
            line width=0.7pt,
            text width=11.0cm,
            align=center,
            font=\footnotesize,
            inner sep=7pt
        },
        flow/.style={
            ->,
            draw=black!65,
            line width=0.8pt
        },
        distinction/.style={
            draw=red!55!black,
            fill=red!4,
            rounded corners=3pt,
            line width=0.7pt,
            text width=11.0cm,
            align=center,
            font=\footnotesize,
            inner sep=7pt
        }
    ]

    \node[process] (manifold) {
        \textbf{Manifold of Sensible Intuition}\\
        Representations given under the forms of space and time
    };

    \node[process, below=of manifold] (synthesis) {
        \textbf{Synthesis of the Manifold}\\
        Combination according to rules of the understanding
    };

    \node[central, below=of synthesis] (apperception) {
        \textbf{Transcendental Unity of Apperception}\\
        The manifold must be capable of belonging to one
        self-consciousness: the unity of the ``I think''
    };

    \node[result, below=of apperception] (categories) {
        \textbf{Categories and Objective Unity}\\
        The manifold is represented under rules that make
        the cognition of objects possible
    };

    \node[result, below=of categories] (experience) {
        \textbf{Unified Possible Experience}\\
        Objects can be represented within one rule-governed
        experience
    };

    \draw[flow] (manifold) -- (synthesis);
    \draw[flow] (synthesis) -- (apperception);
    \draw[flow] (apperception) -- (categories);
    \draw[flow] (categories) -- (experience);

    \node[limit, below=0.85cm of experience] (limit) {
        \textbf{Critical limit of the inference}\\
        The unity required for objective experience does not,
        by itself, establish that the transcendental subject
        is a simple, enduring, substantial soul.
    };

    \node[distinction, below=0.55cm of limit] (ontology) {
        \textbf{Ontological question left open}\\
        What constitutes the identity of a being independently
        of the conditions under which it is represented as
        an object of possible experience?
    };

    \end{tikzpicture}

    \caption{Kant's account of synthesis, transcendental
    apperception, and the unity of possible experience.
    The diagram distinguishes the conditions of objective
    cognition from the ontological constitution of the
    subject or of the objects represented.}

    \label{fig:VS.B01.P01.C08.S10.SS07.P006.DG01}
\end{figure}

% =============================================================================
% END OF DIAGRAM DG01
% =============================================================================